\subsection{特征多项式的性质：迹与行列式}

设 E 是交换域 K 上有限维数为 n 的向量空间，且 u 是 E 的一个自同态。由 III，第 541 页可知，u 的特征多项式具有如下形式：

\[
\chi_ {u} (\mathbf {X}) = \mathbf {X} ^ {n} - \operatorname{Tr} (u) \mathbf {X} ^ {n - 1} + \dots + (- 1) ^ {n} \det (u).\tag{5}
\]

命题10. --- 设 E 是交换域 K 上有限维数为 n 的向量空间，设 u 是 E 的一个自同态，且设 \(\chi_u(X) = \prod_{i=1}^{n}(X - \alpha_i)\) 是其特征多项式的一个线性因式分解

（在 K 的适当扩域中，参见 V，第 21 页）。若 q 是系数在 K 中的多项式，则 \(q(u)\) 的特征多项式由以下公式给出

\[
\chi_ {q (u)} (\mathbf {X}) = \prod_ {i = 1} ^ {n} \left(\mathbf {X} - q \left(\alpha_ {i}\right)\right),\tag{6}
\]

其迹与行列式由下式给出

\begin{enumerate}
\def\labelenumi{(\arabic{enumi})}
\setcounter{enumi}{6}
\tightlist
\item
\end{enumerate}

\[
\operatorname{Tr} (q (u)) = \sum_ {i = 1} ^ {n} q (\alpha_ {i}),\tag{8}
\]

\[
\det (q (u)) = \prod_ {i = 1} ^ {n} q (\alpha_ {i}).
\]

显然，(7) 和 (8) 借助 (5) 由 (6) 得出。为证明公式 (6)，我们可以假设 K 是代数闭的。于是我们取 E 的一组基，使 u 的矩阵 U 关于它是下三角的（VII，第 35 页，命题 8 的推论）；我们将利用以下简单的引理：

引理1. --- 若 B 和 C 是 n 阶下三角矩阵，其对角线分别为 \((\beta_{i})\) 与 \((\gamma_{i})\)，则矩阵 \(B + C\) 与 BC 是下三角矩阵，其对角线分别为 \((\beta_{i} + \gamma_{i})\) 与 \((\beta_{i}\gamma_{i})\)。

设 u 的矩阵 U 是下三角矩阵，其对角线为 \((\alpha_{i})\)，由引理 1 可知 \(q(U)\) 是对角线为 \((q(\alpha_{i}))\) 的下三角矩阵。于是 \(X \cdot I_{n} - q(U)\) 是对角线为 \((X - q(\alpha_{i}))\) 的下三角矩阵，这就证明了 (6)。

推论1. --- 要使 \(q(u)\) 可逆，必要且充分的是 \(q\) 与 \(\chi_u\) 互素。

事实上，说 q 和 \(\chi_{u}\) 互素等价于说它们在 K 的某个代数闭扩张中没有公共根，换句话说（由 (8)）即 \(\det(q(u)) \neq 0\)。

注记1. --- 一个多项式与 \(\chi_{u}\) 互素，当且仅当它与 \(u\) 的最小多项式互素（VII，第 33 页，推论 2）。

推论 2. --- 设 \(r \in K(X)\) 是 \(K\) 上的一个有理函数。则 \(u\) 在 \(r\) 中是可代入的（IV，第 21 页），当且仅当它的每个特征值都是可代入的。在这种情形下，以下公式成立：

\[
\chi_ {r (u)} (\mathrm{X}) = \prod_ {i = 1} ^ {n} \left(\mathrm{X} - r \left(\alpha_ {i}\right)\right), \quad \operatorname{Tr} (r (u)) = \sum_ {i = 1} ^ {n} r \left(\alpha_ {i}\right), \quad \det (r (u)) = \prod_ {i = 1} ^ {n} r \left(\alpha_ {i}\right).
\]

写 r = p/q，其中 p 和 q 是互素的多项式。则 u 在 r 中可代入当且仅当 \(\det(q(u)) \neq 0\)，因此第一个断言由 (8) 得出。根据推论 1，我们可以假设 q 与 \(\chi_{u}\) 互素，于是由贝祖恒等式，存在多项式 g 和 h 使得 \(qg + h\chi_{u} = 1\)。从而由凯莱–哈密顿定理（III，第 541 页）得 \(q(\alpha_{i})g(\alpha_{i}) = 1\) 且 \(q(u)g(u) = 1\)。然后，把公式 (6)、(7) 和 (8) 应用于 \(p(u)g(u) = r(u)\)，即得所述公式。

推论3. --- 对每个整数 \(s \geqslant 0\)，我们有 \(\operatorname{Tr}(u^s) = \sum_{i=1}^{n} \alpha_i^s\)。此公式对 \(s < 0\) 也成立，只要 \(u\) 是可逆的。

这是前述推论的一个特例。

推论4. --- 设域 K 的特征为零；则自同态 u 是幂零的，当且仅当 \(\operatorname{Tr}(u^{s}) = 0\)（对 \(1 \leqslant s \leqslant n\)）。

如果 \(u\) 是幂零的，则诸 \(\alpha_{i}\) 全为零，且 \(\operatorname{Tr}(u^{s}) = 0\) 对所有 \(s > 0\) 成立（推论 3）。反之，如果 \(\operatorname{Tr}(u^{s}) = 0\)（对 \(1 \leqslant s \leqslant n\)），则诸 \(\alpha_{i}\) 全为零，因为 \(K\) 的特征为零（IV，第 72 页，推论），且 \(u\) 是幂零的（VII，第 33 页）。

推论5. --- 设 Y 是一个未定元，令 \(\tilde{u}\) 表示 K(Y)-向量空间 K(Y) \(\otimes_{\mathbf{K}}\) E 的由 \(u\) 经标量扩张（从 K 到系数取自 K 的、关于 Y 的有理函数域 K(Y)）诱导出的自同态。则自同态 Y . 1 - \(\tilde{u}\) 是可逆的。此外，若 \(\chi_u'\) 表示多项式 \(\chi_u\) 的导数，则

\[
\operatorname{Tr} \left(\left(\mathrm{Y}. 1 - \tilde {u}\right) ^ {- 1}\right) = \chi_ {u} ^ {\prime} (\mathrm{Y}) / \chi_ {u} (\mathrm{Y}).
\]

自同态 Y.1 - \(\tilde{u}\) 是可逆的，因为它的行列式是 K(Y) 中的非零元素 \(\chi_u(Y)\)。由此可知 \(\tilde{u}\) 在 K(Y)(X) 的有理函数 \(r(X) = (Y - X)^{-1}\) 中是可代入的。第二个断言现由推论 2 借助关系式得出

\[
\chi_ {u} ^ {\prime} (\mathrm{Y}) / \chi_ {u} (\mathrm{Y}) = \sum_ {i} \left(\mathrm{Y} - \alpha_ {i}\right) ^ {- 1} = \sum_ {i} r \left(\alpha_ {i}\right).
\]

推论 6. --- 设域 K 的特征为零。那么，在形式幂级数环 K{[}{[}T{]}{]} 中，我们有

\[
- \mathrm{T} \frac {d}{d \mathrm{T}} \log \det (1 - \mathrm{Tu}) = \sum_ {m \geqslant 1} \operatorname{Tr} \left(u ^ {m}\right) \mathrm{T} ^ {m}.
\]

让我们首先在有理函数域 \(\mathrm{K}(\mathrm{T})\) 中工作，并令 \(\mathrm{P}(\mathrm{T}) = \det(\mathrm{I}_{n} - \mathrm{T} \cdot U)\)，其中 U 是 u 关于 E 的某一组基的矩阵。则

\[
\mathrm{P} (\mathrm{T}) = \det \left(\mathrm{T} \left(\mathrm{T} ^ {- 1} \cdot I _ {n} - U\right)\right) = \mathrm{T} ^ {n} \chi_ {U} \left(\mathrm{T} ^ {- 1}\right),
\]

所以 \(\mathrm{P}'(\mathrm{T}) / \mathrm{P}(\mathrm{T}) = n / \mathrm{T} - \chi_u'(\mathrm{T}^{-1}) / \mathrm{T}^2\chi_u(\mathrm{T}^{-1})\)。此外，由推论 5 我们有

\[
\chi_ {U} ^ {\prime} \left(\mathrm{T} ^ {- 1}\right) / \mathrm{T} \chi_ {U} \left(\mathrm{T} ^ {- 1}\right) = \operatorname{Tr} \left(\left(\mathrm{T} ^ {- 1} \cdot I _ {n} - U\right) ^ {- 1}\right) / \mathrm{T} = \operatorname{Tr} \left(\left(I _ {n} - \mathrm{T} \cdot U\right) ^ {- 1}\right).
\]

由此可知 \(-\mathrm{TP}^{\prime}(\mathrm{T})/\mathrm{P}(\mathrm{T})=-n+\mathrm{Tr}\left((I_{n}-\mathrm{T}U)^{-1}\right)\) 。现在通过将等式两边展开为形式幂级数即可得到该推论。

注记2. --- 由 IV，第 80 页，推论 1 及公式 (8)，对每个多项式 \(q \in \mathbf{K}[\mathbf{X}]\)，我们有

\[
\det q (u) = \operatorname{res} \left(\chi_ {u}, q\right),\tag{9}
\]

其中 \(\operatorname{res}(\chi_{u}, q)\) 是多项式 \(\chi_{u}\) 与 q 的结式。特别地，若取 \(q = \chi_{u}'\)，我们得到

\[
\det \chi_ {u} ^ {\prime} (u) = (- 1) ^ {n (n - 1) / 2} \mathrm{dis} (\chi_ {u}),\tag{10}
\]

其中 \(\mathrm{dis}(\chi_{u})\) 是多项式 \(\chi_{u}\) 的判别式（IV，第 82 页，公式 (47)）。此外：

推论7. --- 我们有 \(\det\left(\operatorname{Tr}\left(u^{i+j}\right)_{0\leqslant i,j\leqslant n}\right)=\operatorname{dis}\left(\chi_{u}\right)\)。

设 D 为范德蒙德矩阵 \((\alpha_j^{i-1})_{1 \leqslant i, j \leqslant n}\) 。则（III，第 532 页，公式 (29)）：

\[
\det (D) ^ {2} = \prod_ {i <   j} (\alpha_ {i} - \alpha_ {j}) ^ {2} = \operatorname{dis} (\chi_ {u}).
\]

此外，\((i,j)\) 元（\(D\cdot {}^t D\) 的）为 \(\sum_{k}\alpha_{k}^{i + j - 2} = \mathrm{Tr}(u^{i + j - 2})\)，推论随之成立。

